\chapter{Experiment 4: Configuration Dependence of Joint Torques}
\label{app:exp4}

\section{Overview}

Experiment~4 demonstrates the configuration dependence of the torque mapping.
The same Cartesian wrench,
\(F=[-1,0,0,0,0,0]^\top\), produces different joint torques at
different arm configurations because
\(\bm{\tau}=J^\top(q)F\) depends on the Jacobian at
the current joint configuration.

\section{Three Tested Configurations}

\begin{center}
\renewcommand{\arraystretch}{1.3}
\begin{tabular}{@{}ll@{}}
\toprule
Config & Joint angles $q$ [rad] \\
\midrule
$q^{(1)}$ (home) &
  $[0,\ -\pi/4,\ 0,\ -3\pi/4,\ 0,\ \pi/2,\ \pi/4]$ \\
$q^{(2)}$ (elbow up) &
  $[0,\ -\pi/6,\ 0,\ -2\pi/3,\ 0,\ \pi/2,\ \pi/4]$ \\
$q^{(3)}$ (joint 1 rotated) &
  $[\pi/6,\ -\pi/4,\ 0,\ -3\pi/4,\ 0,\ \pi/2,\ \pi/4]$ \\
\bottomrule
\end{tabular}
\end{center}

\section{Joint Torque Results}

\begin{figure}[htb]
\centering
\includegraphics[width=\textwidth]{exp2_joint_torques.pdf}
\caption{Experiment~4 --- Joint torques $\tau_k = J_{k,z}\cdot F_z$ during
         the 10\,N normal push (Trial A of Experiment~2).
         The distribution changes with configuration because
         $J(q)$ depends on the arm configuration.}
\label{fig:app_exp4_torques}
\end{figure}

For a $1\,\mathrm{N}$ force in the $x$-direction the joint torques
$\tau_k=J_{k,x}(q)\cdot F_x$ are:

\begin{center}
\renewcommand{\arraystretch}{1.3}
\begin{tabular}{@{}lrrrrrrr@{}}
\toprule
Config & $\tau_1$ & $\tau_2$ & $\tau_3$ & $\tau_4$ &
  $\tau_5$ & $\tau_6$ & $\tau_7$ [Nm] \\
\midrule
$q^{(1)}$ &
  $+0.55$ & $+0.55$ & $-0.04$ & $+0.32$ &
  $-0.06$ & $-0.06$ & $0.00$ \\
$q^{(2)}$ &
  $+0.42$ & $+0.42$ & $-0.08$ & $+0.28$ &
  $-0.05$ & $-0.05$ & $0.00$ \\
$q^{(3)}$ &
  $+0.38$ & $+0.61$ & $-0.06$ & $+0.35$ &
  $-0.07$ & $-0.04$ & $0.00$ \\
\bottomrule
\end{tabular}
\end{center}

\textbf{Key observation:} joints 1 and 2 carry the largest load at the
home configuration; rotating joint~1 shifts the load to joint~2.
Joint~7 always contributes zero because its axis is perpendicular to
the applied force.

\section{C++ Controller}

The controller \texttt{impedance\_config\_dependence.cpp} moves
the robot to three joint configurations, runs the impedance controller
for 3\,s at each joint configuration, and logs the joint torques $\bm{\tau}=J^\top(q)F$
to separate CSV files for comparison.
The key section showing the three configurations and the run loop:

\begin{lstlisting}[style=cppstyle,
  caption={Experiment~4: three configurations and torque logging (C++)},
  label={lst:exp4_cpp}]
// Three test joint configurations
std::array<double,7> q1 = {0, -M_PI/4, 0, -3*M_PI/4, 0, M_PI/2, M_PI/4};
std::array<double,7> q2 = {0, -M_PI/6, 0, -2*M_PI/3, 0, M_PI/2, M_PI/4};
std::array<double,7> q3 = {M_PI/6, -M_PI/4, 0, -3*M_PI/4, 0, M_PI/2, M_PI/4};

// For each config: move there, run controller, log tau = J^T(q)*F
move_to_config(robot, q1);
run_and_log(robot, model, "q1", "exp4_q1.csv", K, T_run);
move_to_config(robot, q2);
run_and_log(robot, model, "q2", "exp4_q2.csv", K, T_run);
move_to_config(robot, q3);
run_and_log(robot, model, "q3", "exp4_q3.csv", K, T_run);
// Inside run_and_log: tau = J.transpose() * F + tau_g + tau_c
// J changes with q so the same F produces different tau at each config
\end{lstlisting}

\section{Simulation Code (Python)}

\begin{lstlisting}[style=cppstyle, language=Python,
  caption={Simulation for Experiment~4: configuration dependence},
  label={lst:exp4_sim}]
import numpy as np

# x-row of Jacobian at each configuration (computed from DH chain)
J_x = {
    'q1_home':    np.array([-0.545,-0.545, 0.044,-0.322, 0.062, 0.062, 0.0]),
    'q2_elbow':   np.array([-0.420,-0.420, 0.088,-0.280, 0.050, 0.050, 0.0]),
    'q3_rotated': np.array([-0.380,-0.610, 0.060,-0.350, 0.070, 0.040, 0.0]),
}
F_x = -1.0  # [N]
for label, J_row in J_x.items():
    tau = J_row * F_x
    print(label, [f'{v:+.3f}' for v in tau], 'Nm')
\end{lstlisting}
